\documentstyle[%


righttag%


,11pt%


,amssymb%


,russian%               remove in Eng.


,izvru%                 Set izven% in Eng.


]{amsart}





% CHANGE BOTH BMP!!


% J.L.Cie\'sli\'nski


\newcommand{\rf}[1]{(\ref{#1})}


\newcommand{\const}{\operatorname{const}}


\newcommand{\Spin}{\operatorname{Spin}}


\newcommand{\Pin}{\operatorname{Pin}}


\newcommand{\tg}{\operatorname{tg}}


\newcommand{\ctg}{\operatorname{ctg}}


\newcommand{\arctg}{\operatorname{arctg}}


% \newcommand{\arcsin}{\operatorname{arcsin}}


\newcommand{\Cl}{{\cal C}}


\newcommand{\e}{\bold e}





\newcommand{\tts}[1]{}





\theoremstyle{definition}


\theorembodyfont{\rm}


\newtheorem{probl}{\hskip\parindent Проблема}





%\hoffset=-15mm%                  For Eng.


\setcounter{page}{19}


\god{2004}


\nomer{10~(509)} %                        Remove in Eng.


%\nomer{Vol.\,48, No.\,10}%               Set in Eng.


%\str{\pageref{begin}--\pageref{end}}%  Set in Eng.


\udk{514.764}





\author{\it Ю.A.\,Аминов, Я.\,Чешлинский}


% NB:   ^^ remove in Eng.


\title{Изометрические погружения областей пространства


Лобачевского в сферы и евклидовы пространства


и геометрическая интерпретация спектрального параметра}


\thanks{Исследования первого автора частично поддержаны Государственным


фондом фундаментальных исследований Украины (проект \No\,01.07/00132).


Исследования второго автора частично поддержаны


польским комитетом по научным исследованиям, грант KBN \No\,2~P03B~126~22.}





\date{}


%\pagestyle{myheadings}%         Set in Eng.


\pagestyle{plain} %              Remove in Eng.


\begin{document}


%\thispagestyle{plain}%          Set in Eng.


\maketitle


\label{begin}











\section{Введение}





В теории солитонных уравнений (или вполне интегрируемых нелинейных систем) 


\cite{ZMNP} важную роль играет спектральный параметр.


Если в основе интегрируемой системы лежит геометрия


(напр., как в случае уравнения ``синуса-Гордона''


$\varphi,_{xy} = \sin\varphi$, которое описывает


псевдосферические погружения), то


можно ожидать, что спектральный параметр


может быть интерпретирован геометрически.


Действительно, было обнаружено, что в некоторых случаях


геометрическое преобразование Бэклунда зависит от спектрального параметра


(напр., \cite{Sym}).





Напомним две идеи, важные в этом контексте.


Во-первых, использование спектрального параметра дает


естественную конструкцию


локального погружения с помощью так называемой формулы


Сыма--Тайфеля (Sym--Tafel) \cite{Ci-FG}--\cite{S-2}.


Например, начиная со спектральной задачи для уравнения ``синус-Гордона'',


приходим к псевдосферическим поверхностям \cite{S-2}, \cite{Sym}.


Преобразования Бианки--Бэклунда восстанавливаются автоматически.





Во-вторых, Долива и Сантини, рассматривая эволюцию кривых на сфере $S^n$,


получили линейные задачи для ассоциированных нелинейных уравнений.


Они обнаружили, что радиус объемлющей сферы


связан со спектральным параметром \cite{DS0}, \cite{DS2}.


В предельном случае (бесконечный радиус) $S^n \rightarrow E^n$


они вывели формулу Сыма--Тайфеля для эволюции кривых.


Естественно предположить, что спектральный параметр связан с радиусом


объемлющей сферы и в более общей ситуации.





В данной статье представляется полное обсуждение обеих задач


в случае интегрируемых систем, ассоциированных с погружениями


пространства Лобачевского в евклидовы пространства


\cite{Am77}--\cite{Te}.


В разделах \ref{L2}, \ref{Ln} показывается,


что система уравнений Гаусса и Вейнгартена


для погружения $n$-мерного пространства Лобачевского в сферу


$S^{2n-1}$ радиуса $R$ может быть


интерпретирована как спектральная задача


для погружения $n$-мерного пространства Лобачевского


в евклидово пространство $E^{2n-1}$.


Спектральный параметр $\lambda$ есть функция радиуса $R$.


Полученная спектральная задача эквивалентна линейной задаче,


ранее рассмотренной в \cite{ABT} и \cite{FP}.


Используя алгебру Клиффорда


$\Cl(2n)$ и группу $\Spin(2n)$, выводим


формулу Сыма--Тайфеля в предельном случае $R\rightarrow \infty$


(см. раздел 5).


Этот результат был ранее сообщен


(без полного доказательства) в \cite{CA}.





В разделе \ref{obrot} обсуждается некоторый специальный класс


подмногообразий с постоянной отрицательной кривизной,


погруженных в сферу $S^{2n-1}$.


Этот класс является обобщением тора Клиффорда.


Случай $n=2$ оказывается выделенным.


Для $n\neq 2$ получаются


только подмногообразия, порождаемые кривой, подобной трактрисе


(см. рис.~1a)).


Для $n=2$ имеется три класса подмногообразий, включающие


компактные поверхности (хотя и с особенностями). 


\begin{center}


\unitlength=1mm


\begin{picture}(170,97) %% width, heigth, added 5 mm for signature


\put(2,97){\special {em:graph am-ch1.bmp}} %% left X, upper Y, -, /2, +toX


\put(78,0){\mbox Рис.\,1}


\end{picture}


\end{center}














\section{Предварительные замечания} \label{Prel}





Рассмотрим изометрические погружения областей $L^n$


(т.\,e. $n$-мерного пространства постоянной отрицательной кривизны)


в сферу $S^{2n-1}$ с радиусом $R$. 


Пусть $r=r(u^1,\dotsc,u^n)$ --- радиус-вектор погружения, которое 


рассматривается как подмногообразие (возможно с особенностями)


в $E^{2n}$, где $u^1,\dotsc,u^n$ --- координаты кривизны.


Пусть $\nu_1,\dotsc,\nu_{n-1}$, $\nu_n$ --- ортонормированный базис


нормалей к этому подмногообразию. 


Будем предполагать, что один нормальный вектор


параллелен радиус-вектору объемлющей сферы, т.\,е. $\nu_n\equiv r/R$.


В этом случае уравнения Гаусса--Вейнгартена запишутся следующим


образом:


$$


r,_{ij} = \Gamma_{ij}^k r,_{k} + L_{ij}^\alpha \nu_\alpha;\quad


\nu_\alpha,_{i} = - L_{ij}^\alpha g^{jk} r,_{k} + \mu_{\beta\alpha|i} \nu_\beta\ ,


$$


где запятая означает частную производную


($r,_i := \partial r/\partial u^i$ и т.д.),


$\Gamma_{ij}^k$ --- символы Кристоффеля,


$L_{ij}^k$ ($k=1,\dotsc,n$) ---


коэффициенты вторых квадратичных форм


и $\mu_{\beta\alpha|i}$ --- коэффициенты кручения.


По повторяющимся индексам предполагается суммирование.


Можно показать, что в нашем случае нормальный пучок является плоским,


т.\,е. существуют нормальные векторные поля


$\nu_1,\dotsc,\nu_n$ такие, что коэффициенты


кручения тождественно равны нулю \cite{Mo}.





В разделе 5 используем клиффордову алгебру $\Cl(2n)$ и


группу $\Spin(2n)$ для описания вращений в $E^{2n}$.


Клиффордова алгебра $\Cl(m)$ содержит единицу~1 и


порождается элементами


$\e_1,\dotsc,\e_{m}$, которые удовлетворяют соотношениям


$$


\e_j \e_k = - \e_k \e_j \ \ (j\neq k), \quad 


\e_j^2 = 1 \ \ (j=1,\dotsc,m).


$$


Поэтому каждое число Клиффорда (некоторый элемент $\bold v$ из $\Cl(m)$)


может быть представлено в виде


$$


\bold v=v^0 + v^i \e_i + v^{ij} \e_i \e_j + \dotsb +


v^{i_1i_2\dotsc i_k} \e_{i_1} \e_{i_2} \dotsc \e_{i_k} + \dotsb


+ v^{12\dotsc m} \e_1\e_2\dotsc\e_{m} \ ,


$$


где $i < j$, $i_1 < i_2 < \dotsb < i_k$ и т.\,д., и


$v^0,v^i, \dotsc, v^{12\dotsc m}$ --- вещественные коэффициенты.





Обозначим линейное пространство, порождаемое элементами $\e_1,\dotsc,\e_{m}$,


через $V$. Определяя скалярное произведение в $V$ формулой


$(v,w)=(vw + wv)/2$,


отождествляем $V$ с евклидовым пространством $E^{m}$. 


Элементы $\e_1, \dotsc, \e_{m}$ образуют ортонормированный базис


пространства\linebreak $V \simeq E^m$.





Ортогональные преобразования в $E^m$ могут быть


описаны в терминах чисел Клиффорда.


Действительно, пусть $\nu$ --- единичный вектор


(нормаль к некоторой гиперплоскости $\pi$) и $w$ --- произвольный


элемент из $V$.


Легко проверить, что элемент


$w' := - \nu w \nu^{-1}$


также принадлежит $V$ и, более того, $w'$ есть отражение


элемента $w$ относительно гиперплоскости $\pi$ (ортогональной к $\nu$).





Как известно, любое ортогональное преобразование


в $E^m$ есть композиция конечного числа


отражений относительно гиперплоскостей.


Поэтому любое ортогональное преобразование в $E^m$


может быть представлено в виде


$\wt{w} = (-1)^K g_K w g_K^{-1}$,


где $g_K$ --- произведение $K$ единичных векторов.


Заметим, что $g_K$ и $-g_K$ описывают одно и то же преобразование.





Элементы алгебры Клиффорда $\Cl(m)$, которые могут быть разложены


в произведение единичных векторов, 


образуют мультипликативную группу по отношению к умножению Клиффорда.


Эта группа известна как $\Pin(m)$.


Элементы алгебры $\Cl(m)$, представляющиеся в виде


произведения четного числа


единичных векторов, образуют группу $\Spin(m)$.


Ортогональные преобразования,


соответствующие элементам $\Spin(m)$, сохраняют ориентацию.


Группа $\Pin(m)$ есть двойное накрытие группы $O(m)$,


в то время как $\Spin(m)$ есть двойное накрытие группы $SO(m)$.


Поэтому любая кривая $\gamma \subset O(n)$ может


быть поднята на $\Pin(m)$ (двумя способами).


Oбозначим одну из этих поднятых кривых через $S(\gamma)$:


$$


O(n) \supset \gamma(t) \mapsto S(\gamma(t)) \in \Pin(n).


$$


То же касается других подмногообразий в $O(n)$.


Более подробную информацию, относящуюся к


алгебрам Клиффорда, можно найти, например, в \cite{BT}, \cite{Po}.








\section{Локальные погружения $L^2$ в $S^3$} \label{L2}





Начнем с простейшего случая.


Рассмотрим поверхность $F^2 \subset S^3


\subset E^4$, которая локально изометрична $L^2$.


Пусть $u^1$, $u^2$ --- координаты кривизны на $F^2\subset S^3$.


Нормальное пространство подмногообразия $F^2\subset E^4$


имеет размерность~2.


Bыберем одну нормаль, например, $\nu_1$, касательную к сфере


$S^3$, и вторую нормаль в виде $\nu_2 = r/R$.


Коэффициенты кручения $\mu_{12|i}$, $i=1,2$, обращаются в нуль.


Действительно,


$\mu_{12|i} = (\nu_1 \nu_2,_i) = R^{-1} (\nu_1 r,_i) = 0$.


Так как главные направления ортогональны друг другу, можно ввести 


такие ортогональные координаты, что


$ds^2 = g_{11} du_1^2 + g_{22} du_2^2$,


$\mathrm{II}^1 = L_{11}^1 du_1^2 + L_{22}^1 du_2^2$.


Из уравнений Гаусса для поверхности $F^2 \subset S^3$


($K_i = K_e + K_{S^3}$) имеем


$-1 = \frac{L_{11}^1 L_{22}^1}{g_{11} g_{22}} + \frac{1}{R^2}$.


Обозначим


\begin{equation} \label{a2}


a^2 := 1 + {1}/{R^2},


\end{equation}


тогда, очевидно,


$\frac{L_{11}^1 L_{22}^1}{g_{11} g_{22}} = - a^2$.


Поэтому можно ввести такую функцию $\omega$, что 


$\frac{L_{11}^1}{a g_{11}} = \tg\omega$,


$\frac{L_{22}^1}{a g_{22}} = - \ctg \omega$.


Хорошо известно, что на любой поверхности постоянной гауссовой


кривизны $K=-1$ существуют координаты,


в которых метрика может быть записана в виде


$ds^2 = \cos^2\omega du^2 + \sin^2 \omega dv^2$.


Тогда


$L_{11}^1 = a\sin\omega\cos\omega$,


$L_{22}^1 = - a\sin\omega\cos\omega$.


Коэффициенты второй квадратичной формы по отношению ко второй нормали


(ортогональной к сфере $S^3$) легко вычисляются:


$$


L_{11}^2 = (r,_{uu} \nu_2) = R^{-1} (r,_{uu} r) = - \frac{r,_u^2}{R} =


- \frac{\cos^2\omega}{R},\quad


L_{22}^2 = - \frac{r,_v^2}{R} = - \frac{\sin^2\omega}{R}.


$$


Пусть $\tau_1$, $\tau_2$ --- единичные векторы,


касательные к координатным линиям


$$


\tau_1 := \frac{r,_u}{\sqrt{g_{11}}}, \quad 


\tau_2 := \frac{r,_v}{\sqrt{g_{22}}}.


$$


Используя уравнения Гаусса


$r,_{ij} = \Gamma_{ij}^k r,_k + L_{ij}^\alpha \nu_\alpha$,


получим





\begin{gather}


\frac{\partial \tau_1}{\partial u} = \frac{\partial\omega}{\partial v} \tau_2


+ a \nu_1 \sin\omega - \frac{\cos\omega}{R} \nu_2 \ ,\notag\\


%


\label{GW-Gauss}


\frac{\partial \tau_1}{\partial v} = \frac{\partial\omega}{\partial u} \tau_2


\ , \quad \frac{\partial \tau_2}{\partial u} = - 


\frac{\partial\omega}{\partial v} \tau_1 \ ,\\


%


\frac{\partial \tau_2}{\partial v} = - \frac{\partial\omega}{\partial u} \tau_1


- a \nu_1 \cos\omega - \frac{\sin\omega}{R} \nu_2.\notag


\end{gather}


Базис $\nu_1$, $\nu_2$ переносится параллельно в нормальном пучке


(т.\,е. $\mu_{\alpha\beta|i} = 0$).


Используя уравнения Вейнгартена, находим


\begin{equation}


\label{GW-Weingarten}


\begin{split}


\nu_1,_u&=-a\tau_1 \sin\omega \ , \quad \nu_1,_v=- a \tau_2 \cos\omega \ ,\\


%


\nu_2,_u = &- R^{-1} \cos\omega \tau_1 \ , \quad


\nu_2,_v = - R^{-1} \sin\omega \tau_2.


\end{split}


\end{equation}


Уравнения \rf{GW-Gauss} и \rf{GW-Weingarten}


могут быть переписаны в матричной форме


\begin{equation}


\label{prlin}


\begin{split}


\frac{\partial}{\partial u} \begin{pmatrix}


\tau_1 \\ \tau_2 \\ \nu_1\\ \nu_2 \end{pmatrix} &=


\begin{pmatrix} 0 & \omega,_v & a \sin\omega & - R^{-1} \cos\omega \\


- \omega,_v & 0 & 0 & 0 \\ - a \sin\omega & 0 & 0 & 0 \\


R^{-1} \cos\omega & 0 & 0 & 0 \end{pmatrix}


\begin{pmatrix} \tau_1 \\ \tau_2 \\ \nu_1\\ \nu_2 \end{pmatrix},\\


%


\frac{\partial}{\partial v} \begin{pmatrix}


\tau_1 \\ \tau_2 \\ \nu_1\\ \nu_2 \end{pmatrix} &=


\begin{pmatrix} 0 & \omega,_u & 0 & 0 \\


- \omega,_u & 0 & -a \cos\omega & - R^{-1} \sin\omega \\ 


0 & a\cos\omega & 0 & 0 \\


0 & R^{-1} \sin\omega & 0 & 0 \end{pmatrix}


\begin{pmatrix} \tau_1 \\ \tau_2 \\ \nu_1\\ \nu_2 \end{pmatrix}.


\end{split}


\end{equation}


Система уравнений Гаусса--Кодацци (совпадающая с условиями совместности


вышеуказанной системы матричных линейных уравнений)


эквивалентна одному уравнению


$\omega,_{uu} - \omega,_{vv} = \sin\omega \cos\omega$,


которое хорошо известно в теории интегрируемых систем 


как уравнение ``синус-Гордона''. 


Hапомним, что уравнения Кодацци уже разрешены


и их следствием является специальная форма метрики.





Важным обстоятельством является то, что условия совместности не зависят


от $R$. Поэтому имеем линейные уравнения, параметризованные


$R$ (второй параметр $a$ зависит от $R$, см. \rf{a2}).


Такая ситуация типична для интегрируемых систем.


Линейная система, содержащая свободный параметр


(``спектральный параметр''), рассматривается как линейная задача


(или пара Лакса). Присутствие параметра


важно для приложения различных методов теории солитонов,


в частности, преобразований Дарбу--Бэклунда.





Удобно выразить $R$ и $a$ в форме рациональных функций другого параметра


$\lambda$:


\begin{equation}


\label{R-a-lambda}


\frac{1}{R} = \frac{1}{2} \left( \frac{1}{\lambda} - \lambda \right),


\quad a = \frac{1}{2} \left( \frac{1}{\lambda} + \lambda \right).


\end{equation}


Если $a > 0$ и $R > 0$, то $0 < \lambda < 1$.


Допуская отрицательные значения $R$ и $a$, можем распространить формулу


\rf{R-a-lambda} на все действительные значения $\lambda$.


Случай $\lambda=1$ соответствует $R=\infty$ (евклидово


пространство $E^{3}$), а для $\lambda \rightarrow 0$ имеем $R \rightarrow 0$


(единственная точка). Pассмотрим эту проблему более подробно


в разделе~\ref{Symfor}.





Пара Лакса \rf{prlin}, рациональная относительно $\lambda$,


является специальным случаем линейных задач, рассматриваемых в работах


\cite{ABT}, \cite{FP}, \cite{Ci-space}


(более подробно в следующем разделе).








\section{Локальные погружения $L^n$ в $S^{2n-1}$} \label{Ln}





Рассмотрим общий случай погружения $L^n \rightarrow S^{2n-1}


\subset E^{2n}$. По аналогии со случаем погружения


$L^n$ в евклидово пространство \cite{Am77}, \cite{Mo} можно доказать,


что существует


$n$ таких главных направлений с ассоциированными координатами 


$u^1,\dotsc,u^n$, что метрика записывается в виде


$ds^2 = \sum\limits_{i=1}^n H_i^2 (du^i)^2$, 


где $H_1,\dotsc,H_n$ --- функции,


удовлетворяющие условию $H_1^2+\dotsb+H_n^2=1$.


Более того, существует такой базис в нормальном пространстве,


что соответствующие коэффициенты кручения равны нулю.


Нормаль $\nu_n=r/R$ принадлежит этому базису. Действительно,


$\mu_{\alpha n|i} = (\nu_\alpha \nu_n,_i) = R^{-1} (\nu_\alpha r,_i )= 0$.


В дальнейшем будем предполагать,


что $\nu_1,\dotsc,\nu_n$ есть такой базис.


Соответствующие квадратичные формы $II^\alpha$ имеют диагональные матрицы:


$\mathrm{II}^\alpha = \sum\limits_{i=1}^n L_{ii}^\alpha (du^i)^2$. 


Уравнения Кодацци $L_{ij}^p;_k - L_{ik}^p;_j = 0$


(точка с запятой означает ковариантную производную) приобретают вид


\begin{equation}


\label{star}


\frac{\partial}{\partial u^k} \sum_{p=1}^{n-1} \left( L_{ii}^p \right)^2


- \frac{1}{g_{ii}} \frac{\partial g_{ii}}{\partial u^k}


\sum_{p=1}^{n-1} \left( L_{ii}^p \right)^2 - g^{kk}


\frac{\partial g_{ii}}{\partial u^k}


\sum_{p=1}^{n-1} L_{kk}^p L_{ii}^p = 0.


\end{equation}





Запишем уравнения Гаусса


$R_{kiki} = \sum\limits_{p=1}^{n-1} L_{kk}^p L_{ii}^p + \frac{1}{R^2} g_{kk} g_{ii}$.


По предположению пространство $L^n$ имеет кривизну $K=-1$, откуда


следует $R_{kiki} = - g_{kk} g_{ii}$. Итак,


\begin{equation} 


\label{Gauss2}


\sum_{p=1}^{n-1} \frac{L_{kk}^p L_{ii}^p}{g_{kk} g_{ii}} = -1-\frac{1}{R^2}.


\end{equation}


Используем те же обозначения, что и в специальном случае $n=2$, т.\,е.


\begin{equation}


\label{aR}


a = \sqrt{\displaystyle 1 + {1}/{R^2} }.


\end{equation}


Уравнения Кодацци \rf{star} можно переписать в виде


$$ 


\frac{\partial}{\partial u^k} \sum_{p=1}^{n-1} \left( L_{ii}^p \right)^2


- \frac{1}{g_{ii}} \frac{\partial g_{ii}}{\partial u^k}


\sum_{p=1}^{n-1} \left( L_{ii}^p \right)^2 + a^2 g_{ii} 


\frac{\partial g_{ii}}{\partial u^k} = 0.


$$


Обозначение


$\ctg\sigma_i :={\sqrt{\sum\limits_{p=1}^{n-1}\left( 


L_{ii}^p \right)^2}}\Big/{(a g_{ii})}$


приводит к упрощению записи уравнений Кодацци


$\frac{\partial}{\partial u^k} ( \log g_{ii} - \log \sin^2\sigma_i) = 0$.


Выбирая подходящую параметризацию на каждой координатной линии,


получим $g_{ii} = \sin^2 \sigma_i$.  Поэтому


$\sum\limits_{p=1}^{n-1} \left( L_{ii}^p \right)^2 =


a^2 \sin^2\sigma_i \cos^2\sigma_i$,


и $L_{ii}^p$ можно переписать следующим образом:


$L_{ii}^p = a \sin\sigma_i \cos\sigma_i \cos\varphi_i^p$,


где


\begin{equation}


\label{sumcos}


\sum_{p=1}^{n-1} \cos^2\varphi_i^p = 1.


\end{equation}


В этих обозначениях уравнения \rf{Gauss2} записываются в виде


\begin{equation}


\label{Gauss3}


\sum_{p=1}^{n-1} \cos\varphi_i^p \cos\varphi_k^p 


= - \tg\sigma_i \tg\sigma_k.


\end{equation}


В полной аналогии со случаем $L^n$ в $E^{2n-1}$ \cite{Am77}


рассматривается матрица


$$


A := \begin{pmatrix} \sin\sigma_1 & \sin\sigma_2 & \cdots & \sin\sigma_n \\


\phi_1^1 & \phi_2^1 & \dotsb & \phi_n^1 \\


\hdotsfor[1.5]{4} \\


\phi_1^{n-1} & \phi_2^{n-1} & \cdots & \phi_n^{n-1} \end{pmatrix},


$$


где $\phi_j^k := \cos\varphi_j^k \cos\sigma_j$.


Теперь уравнения \rf{sumcos} и \rf{Gauss3} могут быть интерпретированы


как условия ортогональности матрицы $A$: $A^\top A = I$.


Отсюда


$\sum\limits_{i=1}^n \sin^2\sigma_i = 1$.


Введем обозначения


$H_i:=\sin\sigma_i$,


$\tau_i := \frac{r,_i}{\sqrt{g_{ii}}}$,


$\beta_{ij} := \frac{1}{H_i} \frac{\partial H_j}{\partial u^i}$ 


($i\neq j$).


Тогда уравнения Гаусса--Вейнгартена \rf{GW-Gauss}, \rf{GW-Weingarten}


примут вид


\begin{gather*}


\tau_i,_i = - \sum_{k=1}^n \beta_{ki} \tau_k + a \cos\sigma_i 


\sum_{p=1}^{n-1} \cos\varphi_i^p \nu_p


- \frac{\sigma_i}{R} \nu_n \ , \quad


%


\tau_i,_j = \beta_{ij} \tau_j \quad (i\neq j),\\


%


\nu_{\alpha},_i = - a \cos\sigma_i \cos\varphi_i^\alpha \tau_i \ , \quad


\nu_n,_i = \frac{\sin\sigma_i}{R} \tau_i \ ,


\end{gather*}


или в матричной форме


\begin{equation}


\label{problinPhi}


\frac{\partial \Phi}{\partial u^i} = \Omega_i \Phi \ ,


\end{equation}


$$


\Omega_i :=


\begin{pmatrix}


0 & \dotsb & \beta_{1i} & \dotsb & 0 & 0 & \dotsb & 0 & 0 \\


\vdots & \dotsb & \vdots & \dotsb & \vdots & \vdots & \dotsb & \vdots &


\vdots \\


- \beta_{1i} & \dotsb & 0 & \dotsb & - \beta_{ni} & a\phi_i^1 &


\dotsb & a\phi_i^{n-1} & - R^{-1} \sin\sigma_i \\


\vdots & \dotsb & \vdots & \dotsb & \vdots & \vdots & \dotsb & \vdots &


\vdots \\


0 & \dotsb & \beta_{ni} & \dotsb & 0 & 0 & \dotsb & 0 & 0 \\


0 & \dotsb & - a\phi_i^1 & \dotsb & 0 & 0 & \dotsb & 0 & 0 \\


\vdots & \dotsb & \vdots & \dotsb & \vdots& \vdots & \dotsb & \vdots &


\vdots \\


0 & \dotsb & - a\phi_i^{n-1} & \dotsb & 0 & 0 & \dotsb & 0 & 0 \\


0 & \dotsb & R^{-1} \sin\sigma_i & \dotsb & 0 & 0 & \dotsb & 0 & 0 


\end{pmatrix},


$$


где $\Phi:= (\tau_1, \tau_2, \dotsc, \tau_n, \nu_1, \dotsc, \nu_n)^\top$ 


(через ${}^\top$ обозначена транспозиция)


и кососимметрические матрицы $\Omega_i$ могут иметь


ненулевые коэффициенты только на $i$-й строке и в $i$-м столбце.


Принимая во внимание равенство \rf{aR} и параметризацию


\rf{R-a-lambda} в матрицах $\Omega_i$,


получим линейную (спектральную) задачу,


которая с точностью до изменения базиса эквивалентна


линейной задаче, рассматриваемой в 


\cite{ABT} (см. также \cite{Ci-space}, \cite{FP}):


$$


\frac{\partial\psi}{\partial x_j} = \frac{\lambda}{2} A_j \psi


+ \frac{1}{2\lambda} B_j \psi + C_j \psi 


$$


(надо обратить внимание на ошибку в формуле (1.17b) в \cite{ABT}),


где 


$$


A_j = \begin{pmatrix} 0 & a_j \\ a_j^\top & 0 \end{pmatrix}, \quad 


B_j = \begin{pmatrix} 0 & \delta a_j \\ a_j^\top \delta & 0 \end{pmatrix}, \quad


C_j = \begin{pmatrix} 0 & 0 \\ 0 & \gamma_j \end{pmatrix},


$$


$a_j$ --- матрица размера $n\times n$, все элементы которой равны


нулю, за исключением $j$-го столбца, имеющего вид $a_{1j},\dotsc,a_{nj}$,


$\delta = \mathrm{diag}(1,-1,\dotsc,-1)$,


а $\gamma_j$ --- кососимметрическая матрица размера


$n\times n$, все элементы которой равны нулю, за исключением


$j$-го столбца, имеющего вид


$(\beta_{1j},\dotsc,\beta_{nj})^\top$, и $j$-й


строки, имеющей вид $(-\beta_{1j},\dotsc,-\beta_{nj})$.


Функции $\psi$ принимают значения в $R^{2n}$.


Чтобы сравнить эту спектральную задачу со спектральной задачей


\rf{problinPhi}, надо приравнять


$a_{1k} = \sin\sigma_k$, $a_{jk} = \phi_{j-1}^k$,


где $j=2,\dotsc,n$ и $k=1,\dotsc,n$.


Система уравнений Гаусса--Кодацци--Риччи,


известная в этом случае как система, обобщающая


уравнение ``синус-Гордона'' или систему ``LE''


(\cite{ABT}, \cite{CA}, \cite{Am88}, \cite{Am98}), 


не зависит от $R$. Таким образом, доказана 





\begin{thm}


Система уравнений Гаусса--Вейнгартена


для локальных погружений пространства Лобачевского $L^n$


в сферу $S^{2n-1}$ радиуса $R$


тождественна линейной задаче $($со спектральным параметром$)$, связанной с


локальными погружениями $L^n$ в


$E^{2n-1}$.


Спектральный параметр $\lambda$ может быть выражен через


$R$, а именно


$$


\frac{1}{R} = \frac{1}{2} \left( \frac{1}{\lambda} - \lambda \right).


$$


\end{thm}





Мы вывели рациональное соотношение между $R$ и $\lambda$ для


$\lambda \in (0,1)$ (сравните с окончанием раздела \ref{L2}).


Однако условия совместности в точности те же самые для любого $\lambda$.


Рассмотрим в деталях эту проблему в следующем разделе.








\section{Геометрический вывод формулы Сыма--Тайфеля}


\label{Symfor}





Хорошо известно, что изометрические погружения


областей пространства Лобачевского $L^2$ в $E^3$ могут быть


в точности описаны так называемой формулой Сыма--Тайфеля


\begin{equation}


\label{Sym}


F = \Psi^{-1} \Psi,_\zeta \ ,


\end{equation}


где $\Psi=\Psi(u^1,u^2,\zeta)$ есть $SU(2)$-значное


решение линейной системы


\begin{equation}


\label{SGproblin}


\Psi,_1 = \begin{pmatrix} i\zeta & - \frac{1}{2} \varphi,_1 \\ \frac{1}{2} 


\varphi,_1 & - i \zeta \end{pmatrix} \Psi \ , \quad


\Psi,_2 = \frac{1}{4i\zeta} \begin{pmatrix} \cos\varphi & \sin\varphi \\ 


-\sin\varphi & -\cos\varphi \end{pmatrix} \Psi \ ,


\end{equation}


а $\zeta$ --- спектральный параметр.


Эта линейная система известна как пара Лакса для уравнения


``синус-Гордона'' $ \varphi,_{12} = \sin{\varphi}$


и имеет ненулевые решения тогда и только тогда, когда $\varphi$ удовлетворяет


уравнению ``синус-Гордона''. 





\begin{thm}[{\series{m}\shape{n}\selectfont \cite{Sym}, \cite{S-2}}]


Если $\Psi$ удовлетворяет системе $\rf{SGproblin}$,


то $su(2)$-значная функция


$F=\Psi^{-1}\Psi,_{\zeta}$ определяет $\zeta$-семейство псевдосферических


погружений в $E^3$.


А именно, $F =\sum\limits_{k=1}^3 \frac{1}{2i} \sigma_k f_k$


$($где $\sigma_k$ --- матрицы Паули$)$ и $f:=(f_1,f_2,f_3)\in E^3$


есть радиус-вектор


псевдосферической поверхности с гауссовой кривизной 


$K=-\zeta^2$. Более того, координаты $u^1$, $u^2$


являются асимптотическими.


\end{thm}





Подход Сыма недавно был применен к общему случаю


изометрических погружений $L^n$ в $E^{2n-1}$ \cite{Ci-space}.


Ниже покажем, как вывести в этом случае аналог формулы


\rf{Sym} из чисто геометрических рассмотрений.


Решение $\Phi$ спектральной задачи \rf{problinPhi} есть функция


со значениями в матричной группе Ли $SO(2n)$


(при условии, что начальные условия принимают значения в этой группе).


Действительно, из $(\det\Phi),_j = \mathrm{Tr}\,\Omega_j \det\Phi$


и $\mathrm{Tr}\,\Omega_j = 0$ имеем $\det\Phi = \const$.


Строки матрицы $\Phi$ можно интерпретировать как ортогональный


базис, ассоциированный с рассматриваемым $\lambda$-семейством погружений.


Без ограничения общности можем выбрать начальные условия так,


что в заданной точке $u_0 :=(u_0^1,\dotsc,u_0^n)$ матрица $\Phi$


является единичной матрицей: $\Phi(u_0,\lambda)=I$.


Тогда $\Phi=\Phi(u,\lambda)$ представляет поворот, переводящий


канонический базис в $u_0$ в базис в $u$.


Для того чтобы вывести формулу Сыма--Тайфеля,


перепишем линейную задачу \rf{problinPhi}


в терминах чисел Клиффорда \cite{Ci-space}.


Достаточно отождествить образующие алгебры Ли $so(2n)$


с образующими алгебры Ли группы $\Spin(2n)$, а именно,


$f_{\mu \nu} \ \leftrightarrow \ 2^{-1} \e_\mu \e_\nu$,


где $f_{\mu \nu}$ --- кососимметрические $2n\times 2n$ матрицы с


коэффициентами


$f_{\mu \nu \alpha \beta}: = \delta_{\mu \alpha} \delta_{\nu \beta} -


\delta_{\mu \beta} \delta_{\nu \alpha}$


(другими словами, $f_{\mu\nu}$ имеют только два ненулевых элемента:


1~на пересечении $\mu$-й строки и $\nu$-го столбца и 


$-1$~на пересечении $\nu$-й строки и 


$\mu$-го столбца).


Tакже заменим $a$ и $R$ на $\lambda$ согласно формулам


\rf{R-a-lambda}.


Тогда спектральная задача \rf{problinPhi} примет вид


\begin{equation}


\label{problin-Cliff} 


\begin{split}


&\Psi,_j = U_j \Psi \ , \\


%


U_j := \frac{1}{2} \e_j \bigg( - \sum_{k=1}^{n} \beta_{kj} \e_k +


&\frac{\lambda^2+1}{2\lambda} \sum_{i=1}^{n-1} \phi_j^i \e_{n+i} + 


\frac{\lambda^2-1}{2\lambda} \sin\sigma_j \e_{2n} 


\bigg),


\end{split}


\end{equation}


где $\Psi$ --- элемент $\Spin(2n)$, который соответствует


вращению, определенному $SO(2n)$-значной функцией $\Phi$.


Используя обозначения,


введенные в конце раздела~\ref{Prel}, можем кратко записать


$ \Psi = S (\Phi)$.





\begin{thm}[{\series{m}\shape{n}\selectfont \cite{Ci-space}}]


\label{Cispace}


Формула Сыма--Тайфеля $r = \Psi^{-1} \Psi,_\lambda|_{\lambda=1}$,


где $\Psi$ удовлетворяет соотношениям $\rf{problin-Cliff}$,


приводит к погружению пространства Лобачевского $L^n$


в евклидово пространство


размерности $2n-1$, порождаемое элементами


$\e_k \e_{2n}$ $(k=1,2,\dotsc,2n-1)$.


\end{thm}





Доказательство теоремы 3 проводится непосредственно в \cite{Ci-space}.


Заметим, что в нашем случае формула Сыма--Тайфеля приводит к погружениям


$L^n \subset E^{2n-1}$ 


только для $\lambda=1$, а не для всего $\lambda$-семейства погружений,


как в случае уравнения ``синус-Гордона''. 





Теперь представим геометрический вывод формулы Сыма--Тайфеля.


Рассмотрим погружение $L^n$ в $E^{2n}$.


Напомним, что ассоциированный базис


$(\tau_1,\dotsc,\tau_n$, $\nu_1,\dotsc,\nu_n)^\top$ обозначается через $\Phi$,


$E^{2n}$ отождествлено с пространством $V$, порождающим алгебру $\Cl(2n)$


(см. раздел~\ref{Prel}), а вращение,


задаваемое $SO(2n)$-значной функцией $\Phi$, представлено функцией


$S(\Phi) = \Psi(u,\lambda)\in \Spin(2n)$.


Для заданного элемента $w \in E^{2n}$ будем обозначать


через $\wh{w}$ элемент из $V$, соответствующий $w$


при естественном изоморфизме $V \simeq E^{2n}$. Базис


$(\wh{\tau}_1, \wh{\tau}_2, \dotsc, \wh{\tau}_n$, $\wh{\nu}_1, \dotsc, 


\wh{\nu}_{n-1}, \wh{r}/R)$ 


получается из канонического базиса


$\e_1,\dotsc,\e_{2n}$ следующим образом:


$\wh{\tau}_1 = \Psi^{-1} \e_1 \Psi, \dotsc,


\wh{\tau}_n = \Psi^{-1} \e_n \Psi$,


$\wh{\nu}_1 = \Psi^{-1} \e_{n+1} \Psi, \dotsc, 


\frac{\wh{r}}{R} = \Psi^{-1} \e_{2n} \Psi$.


В дальнейшем используется только последнее уравнение,


с помощью которого получается радиус-вектор рассматриваемого погружения


$L^n \subset S^{2n-1}$:


\begin{equation}


\label{bfr}


\wh{r} = \frac{2\lambda}{1-\lambda^2} \Psi^{-1} \e_{2n} \Psi.


\end{equation}


Формула \rf{bfr} является следствием геометрии погружения


в предположении, что $R > 0$ и $0 < \lambda < 1$.


Возможно формальное применение также и для


$\lambda > 1$.


Тогда из формулы \rf{aR} следует, что $R$ отрицательно.


Объясним геометрический смысл этого случая.


Из элементарных свойств образующих алгебры Клиффорда имеем


(для $k=1,\dotsc,n$) $\e_k \e_\mu = \e_{2n} \e_k \e_\mu \e_{2n}$


($\mu\neq 2n$) и $\e_k \e_{2n} = - \e_{2n} \e_k \e_{2n} \e_{2n}$.


Это приводит к следующей симметрии


спектральной задачи \rf{problin-Cliff}:


$U_k (\tfrac{1}{\lambda}) = \e_{2n} U_k (\lambda) \e_{2n}$.


Поэтому можно ограничиться решениями $\Psi$, удовлетворяющими


условию


\begin{equation}


\label{inv}


\Psi (\tfrac{1}{\lambda}) = \e_{2n} \Psi (\lambda) \e_{2n}.


\end{equation}


Действительно, если вышеприведенное равенство


имеет место при произвольном выборе точки


$(u^1,\dotsc,u^n)$ (начальные условия),


то можно легко доказать (дифференцируя обе части равенства по


$u^k$ и используя соотношения \rf{problin-Cliff}), что оно справедливо


и в окрестности этой точки.


Из соотношений \rf{bfr} и \rf{inv} получаем


\begin{equation}


\label{reflect}


{\wh r} (1/\lambda) = - \e_{2n} {\wh r} (\lambda) \e_{2n}.


\end{equation}


Геометрически это означает, что ${\wh r} (1/\lambda)$ --- отражение


${\wh r}(\lambda)$ относительно гиперплоскости, ортогональной $\e_{2n}$.


Формула \rf{reflect} может рассматриваться как геометрическое определение


${\wh r} (\lambda)$ для $\lambda > 1$.


Другая симметрия спектральной задачи \rf{problin-Cliff}, а именно


$U_k (-\lambda) = \e U_k (\lambda) \e^{-1}$,


где $\e : = \e_{n+1} \e_{n+2} \dotsc \e_{2n}$ ($\e^{-1} = \e_{2n}\dotsc 


\e_{n+1}$), приводит к геометрическому определению ${\wh r} (\lambda)$ для


$\lambda < 0$.


Действительно, выбирая подходящие начальные условия, можем ограничиться


решениями $\Psi$, для которых


$\Psi (- \lambda) = \e \Psi (\lambda) \e^{-1}$.


Тогда


${\wh r} ( - \lambda) = (-1)^n \e {\wh r} (\lambda) \e^{-1}$,


т.\,е. чтобы получить ${\wh r} ( -\lambda)$, нужно произвести


$n$ отражений ${\wh r} (\lambda)$


относительно гиперплоскостей, ортогональных


соответственно $\e_{n+1},\dotsc,\e_{2n}$.


Предельный случай $\lambda=1$ формально соответствует $R=\pm \infty$.


Чтобы получить точное геометрическое значение


предела, рассмотрим формулу \rf{reflect}.


Если $\lambda \rightarrow 1-$, то


${\wh r} (\lambda)$ превращается в погружение в сферу радиуса


$R \rightarrow \infty$, и можно пользоваться формулой \rf{reflect}, чтобы получить


значение правостороннего предела ${\wh r} (\lambda)$ при $\lambda=1$.


В предельном случае погружения ${\wh r} (1+)$ и ${\wh r} (1-)$


симметричны по отношению к гиперплоскости, ортогональной $\e_{2n}$.


Радиус-вектор ${\wh r}$ недостаточно удобен, т.\,к. для


$R \rightarrow \infty$ он стремится к бесконечности.


Следуя идее Доливы и Сантини


(которые использовали ее для случая эволюции кривых, \cite{DS2}),


рассмотрим


$\wh{F} := {\wh r} - R \e_{2n}$,


что просто обозначает изменение базиса.


Теперь подмногообразие описывается


по отношению к фиксированной точке сферы (северный полюс при


$R>0$ и южный полюс при $R<0$), а не к центру сферы.


Тогда


\begin{equation}


\label{hatF1}


\wh{F} = R \Psi^{-1} \e_{2n} \Psi - R \e_{2n} = 


\frac{2\lambda}{1 - \lambda^2} 


( \Psi^{-1} \e_{2n} \Psi - \e_{2n}).


\end{equation}


Можно показать, что


$\wh{F} (1/\lambda) = - \e_{2n} \wh{F} (\lambda) \e_{2n}$.


Это означает, что $\wh{F} (1/\lambda)$ и $\wh{F} (\lambda)$


имеют равные компоненты, ортогональные полярной оси.


Их параллельные компоненты имеют противоположный знак.


$\wh{F}(1-)$ приводит к погружению в сферу бесконечного радиуса


и может быть отождествлено


с погружением в касательное пространство к сфере


в северном полюсе (по крайней мере в случае,


когда рассматривается погружение достаточно малой


области; в случае глобальных погружений вопрос сложнее).


Правосторонний предел в правой части приводит


к тому же погружению, отраженному относительно


гиперплоскости, ортогональной оси $\e_{2n}$.


Итак, векторы $\wh{F} (1+)$


и $\wh{F} (1-)$ принадлежат касательным пространствам


в антиподальных точках.


Рассматривая их как элементы линейного пространства, порождаемого


векторами


$\e_1,\dotsc,\e_{2n-1}$, имеем $\wh{F} (1+) = \wh{F} (1-)$,


хотя эти векторы расположены в различных точках.


В дальнейшем вычислим этот предел, предполагая, что решение $\Psi$


спектральной задачи \rf{problin-Cliff} есть


дифференцируемая функция $\lambda$ в окрестности $\lambda=1$.


Чтобы вычислить предел при $\lambda\rightarrow 1$ в формуле \rf{hatF1},


используем правило Лопиталя.


Дифференцируя тождество $\Psi^{-1} \Psi = \bold 1$,


получим $(\Psi^{-1})'= - \Psi^{-1} \Psi' \Psi^{-1}$ 


(где штрих обозначает производную по $\lambda$).


Следовательно,


\begin{equation}


\label{F}


F := \lim_{\lambda \rightarrow 1} \wh{F}(\lambda) = \Psi(1)^{-1} 


\Psi'(1) \Psi(1)^{-1} \e_{2n} \Psi(1) - \Psi(1)^{-1} \e_{2n} 


\Psi'(1).


\end{equation}


Из \rf{inv} следует


\begin{equation}


\label{comm}


\Psi(1) \e_{2n} = \e_{2n} \Psi(1) \ , \quad \e_{2n} \Psi'(1) =


- \Psi'(1) \e_{2n}.


\end{equation}


Используя \rf{comm}, получим из \rf{F} формулу Сыма--Тайфеля


$$


F = 2 \Psi(1)^{-1} \Psi'(1) \e_{2n} = - 2 \e_{2n} 


\Psi^{-1} \frac{\partial \Psi}{\partial \lambda} \bigg|


_{\lambda=1}.


$$


Множитель при $\e_{2n}$ есть проекция


$\Psi^{-1} \Psi,_{\lambda}|_{\lambda=1}$ на $E^{2n-1}$.


Действительно, $\Psi(1)^{-1} \Psi'(1)$ принимает значения


в алгебре Ли группы $\Spin(2n)$ и согласно \rf{comm} антикоммутирует с


$\e_{2n}$. Это означает, что


$ \Psi(1)^{-1} \Psi'(1) = - \frac{1}{2} \e_{2n} \bold f$,


где $\bold f$ принимает значения в подпространстве $V$, порождаемом


элементами $\e_1,\dotsc,\e_{2n-1}$.


Предельный процесс оставляет уравнения Гаусса--Кодацци--Риччи без изменения,


так что результирующая функция $\bold f$ представляет


погружение $L^n$ в $E^{2n-1}$


(это также следует из теоремы~\ref{Cispace}).


Формула \rf{Sym} упрощает явную конструкцию погружений.


Действительно, допустим, что уже нашли базис касательных и нормалей


в каждой точке погружения, тогда (без формулы Сыма) интегрированием 


можно восстановить погружение. 


Формула Сыма--Тайфеля позволяет заменить интегрирование по параметрам


$u^1,\dotsc,u^n$ дифференцированием по параметру $\lambda$.


Конечно, упрощение такого рода возможно только в случаях,


связанных с интегрируемыми системами 


(существование спектрального параметра играет ключевую роль).


Более того, формула Сыма--Тайфеля оказывается весьма полезной


в приложениях методов теории солитонов к дифференциальной геометрии


(напр., \cite{Sym}, \cite{Ci-FG}, \cite{Bob}, \cite{BP}).








\section{Обобщенный тор Клиффорда постоянной отрицательной кривизны}


\label{obrot}





Обозначим через $e_1,\dotsc,e_{2n}$


ортонормированный базис в евклидовом пространстве $E^{2n}$.


Тор Клиффорда есть $n$-мерное плоское подмногообразие ($K=0$) в $S^{2n-1}


\subset E^{2n}$.


Радиус-вектор $r_C = r_C (u_1,\dotsc,u_n)$ имеет вид


\begin{equation}


\label{C-torus}


r_C = \sum_{k=1}^n c_k \left( e_{2k-1} \cos u_k + e_{2k} \sin u_k \right),


\end{equation}


где $c_1,\dotsc,c_n$ --- такие постоянные, что $c_1^2 + c_2^2 + \dotsb


+ c_n^2 = R^2$ (чтобы $r_C$ было погружением в сферу


$S^{2n-1}$ радиуса $R$).


Найдем класс погружений областей $L^n$


в $S^{2n-1}$, которые имеют вид, близкий к торам Клиффорда.


Допустим, что коэффициенты $c_k$ зависят только от $u_1$.


А именно, рассмотрим погружения $L^n \subset S^{2n-1} \subset E^{2n}$


такие, что радиус-вектор $r=r(u_1,u_2,\dotsc,u_n)$ имеет вид


\begin{gather}


r = R (e_1 \cos u_1 + e_2 \sin u_1) \cos\theta + 


R \rho \sin\theta \ , \quad 


\label{gen-Cliff}


\theta=\theta(u_1) \ ,\\


%


\rho = a_2 (e_3 \cos u_2 + e_4 \sin u_2) + \dotsb


+ a_n (e_{2n-1} \cos u_n + e_{2n} \sin u_n)\notag


\end{gather}


и $a_k$ --- постоянные, удовлетворяющие условию


$\sum\limits_{k=2}^n a_i^2 = 1$.


Очевидно, 


$$


r^2 = R^2 \cos^2 \theta + R^2 \sin^2\theta \sum\limits_{k=2}^n a_i^2 = R^2,


$$


поэтому $r$ определяет подмногообразие в сфере $S^{2n-1}$ радиуса $R$.


Касательные векторы имеют вид


\begin{gather*}


r,_1 = R \cos\theta (- e_1 \sin u_1 + e_2 \cos u_1 + \theta'\rho) - R \theta' 


\sin\theta (e_1 \cos u_1 + e_2 \sin u_1)\ , \\


r,_i = R a_i \sin\theta (-e_{2i-1} \sin u_i + e_{2i} \cos u_i) \ ,


\end{gather*}


где $i=2,\dotsc,n$ и $\theta':= \partial \theta/\partial u_1$. 


Поэтому


$g_{11} = R^2 ( \cos^2 \theta + (\theta')^2 )$,


$g_{ii} = R^2 a_i ^2 \sin^2 \theta$ ($i=2,\dotsc,n$),


$g_{ij} = 0$ ($i\neq j$)


и метрика запишется следующим образом:


$$


ds^2 = R^2 (\cos^2\theta + (\theta')^2) du_1^2 + R^2 \sin^2\theta \sum_{i=2}^n


a_i^2 du_i^2.


$$


В дальнейшем обозначим $H_k := \sqrt{g_{kk}}$. Заметим, что


в нашем случае $H_k$ есть функции от одной переменной $u_1$.


Теперь выведем уравнения для $\theta$, которые вытекают из того условия, что


рассматриваемое подмногообразие имеет кривизну $K=-1$.


В этом случае имеем сильное ограничение на тензор Римана


$R_{ijij} = - g_{ii} g_{jj}$ ($i\neq j$),


$R_{ijik} = 0$ ($i\neq j\neq k\neq i$).


Напомним, что


$$


R_{ijij} = - H_i H_j \left( \frac{\partial}{\partial u_i} \left( 


\frac{1}{H_i} \frac{\partial H_j}{\partial u_i} \right) + 


\frac{\partial}{\partial u_j} \left( \frac{1}{H_j} 


\frac{\partial H_i}{\partial u_j} \right) +


\sum_{k\neq i,j} \frac{1}{H_k^2} \frac{\partial H_i}{\partial u_k}


\frac{\partial H_j}{\partial u_k} \right).


$$


Поэтому в нашем случае уравнение


$R_{1i1i}=-g_{11} g_{ii}$ можно переписать в виде


$$


\frac{\partial}{\partial u_1} \left( \frac{1}{H_1} \frac{\partial H_i}{\partial


u_1} \right) = H_1 H_i \ ,


$$


и далее в терминах $\theta$


\begin{equation}


\label{eqth1}


\frac{d}{du^1}


\left( \frac{\theta' \cos\theta}{\sqrt{\cos^2\theta+ \theta'{}^2}} \right)


= R^2 \sin\theta \sqrt{\cos^2 \theta + \theta'{}^2}.


\end{equation}


Интересно отметить, что это уравнение не зависит от $n$


и не содержит постоянных $a_i$. В случае $n=2$ других уравнений нет.





В случае $n>2$ появляются еще два дополнительных уравнения.


Действительно, уравнения $R_{ijij} = - g_{ii}


g_{jj}$ ($i\geq 2, j\geq 2$) можно переписать


$$


\frac{1}{H_1^2} \frac{\partial H_i}{\partial u_1}


\frac{\partial H_j}{\partial u_1} = H_i H_j \ ,


$$


где учтено, что $H_k = H_k (u_1)$ для $k=1,\dotsc,n$.


Подставляя выражения для $H_k$, получим


\begin{equation}


\label{eqth2}


\frac{\theta'{}^2}{\theta'{}^2 + \cos^2\theta} = R^2 \tg^2 \theta.


\end{equation}


Оставшиеся уравнения ($R_{ijik}=0$) имеют вид


$$


\frac{\partial^2 H_i}{\partial u_j \partial u_k} -


\frac{1}{H_k} \frac{\partial H_i}{\partial u_k} 


\frac{\partial H_k}{\partial u_j} - \frac{1}{H_j}


\frac{\partial H_i}{\partial u_j} 


\frac{\partial H_j}{\partial u_k} = 0 


$$ 


и автоматически удовлетворяются, т.\,к. $H_k = H_k (u_1)$.


Следовательно, имеем одно уравнение \rf{eqth1} для $n=2$ и два уравнения


\rf{eqth1}, \rf{eqth2} для $n>2$.


Для обыкновенного дифференциального уравнения второго


порядка \rf{eqth1} можно найти первый интеграл.


Действительно, обозначим


\begin{equation}


\label{W}


W := \frac{\theta'{}^2}{\cos^2\theta+ \theta'{}^2}.


\end{equation}


Очевидно, $0\leq W \leq 1$.


Тогда уравнение \rf{eqth1} приводится к виду


$(\sqrt{W} \cos\theta)' = \frac{R^2 \theta' \sin\theta }{\sqrt{W}}$,


что может быть переписано следующим образом:


$(\sqrt{W} \cos\theta),_\theta = \frac{R^2 \sin\theta}{\sqrt{W}}$.


Умножая обе стороны на $2\sqrt{W}\cos\theta$, имеем


$(W \cos^2\theta),_\theta = 2 R^2 \sin\theta \cos\theta$, 


и интегрируя, получаем первый интеграл для уравнения \rf{eqth1}


\begin{equation} \label{firstint}


W \cos^2\theta = R^2 \sin^2\theta + C \ ,


\end{equation}


где $C$ --- постоянная. Так как $0\leq W \leq 1$, имеем


$- R^2 < C < 1$.


В случае $n>2$ должно удовлетворяться также уравнение \rf{eqth2},


которое переписывается как $W = R^2 \tg^2\theta$.


Подставляя \rf{eqth2} в \rf{firstint}, получим $C=0$.


В случае $n=2$ постоянная $C$ может


принимать и ненулевые значения.


Подставляя \rf{W} в \rf{firstint}, получаем


\begin{equation} \label{quadr}


u_1 = \pm \int \sqrt{\frac{\cos^2\theta-R^2\sin^2\theta-C}{\cos^2\theta


(R^2\sin^2\theta+C)}} d\theta \ .


\end{equation}








{\it Случай $n>2$}.


В этом случае необходимо $C=0$.


Чтобы вычислить интеграл \rf{quadr}, сделаем подстановку


$t = \sqrt{1-(1+R^2)\sin^2\theta}$.


Тогда, выбирая постоянную интегрирования подходящим образом, имеем


$$


|u_1| = \frac{1}{R} 


\log \frac{\sqrt{1+R^2}|\sin\theta|}{1+\sqrt{1-(1+R^2)\sin^2\theta}}


+\arctg \frac{\sqrt{1-(1+R^2)\sin^2\theta}}{R}.


$$


Очевидно, $\sin\theta < \frac{1}{\sqrt{1+R^2}}$.


Следовательно, $0 < |\theta| < \arctg \frac{1}{R}$.


Более того, для $\theta \rightarrow 0$ имеем


$u_1 \rightarrow \pm \infty$.


Для $\theta=\theta_0$ такого, что $\sin\theta_0=\frac{1}{\sqrt{1+R^2}}$,


имеем $\frac{d u_1}{d \theta} = 0$ (см. рис.~1a)).


Следовательно, график функции $\theta=\theta(u_1)$


напоминает трактрису.


\medskip








{\it Случай $n=2$ и $C<0$}.


При $n=2$ постоянная $C$ может принимать ненулевые значения.


Если $C < 0$, то обозначим $C = - \mu^2$, где $0 < \mu < R$.


Тогда


\begin{equation} \label{poch}


\left( \frac{d u_1}{d \theta} \right)^2 = 


\frac{1+\mu^2 - (1+R^2) \sin^2\theta}{(1-\sin^2\theta)(R^2\sin^2\theta-\mu^2)}.


\end{equation}


Правая часть обращается в нуль для таких $\theta_1$, что


$$


\sin\theta_1 = \sqrt{\frac{1+\mu^2}{1+R^2}} 


$$


(заметим, что $\theta_1$ всегда существует).


Если построить в соответствующих точках график


функции $\theta=\theta(u_1)$, то (рис.~1b)) касательные прямые параллельны


$\theta$-оси. Кривые, порождаемые этими точками при вращении вокруг


$u_1$-оси, являются сингулярными кривыми на многообразии.


Для $\theta_2$ такого, что $\sin\theta_2=\mu/R$


(заметим, что $\theta_2$ также всегда существует), 


знаменатель в правой части уравнения \rf{poch}


может обратиться в нуль.


В соответствующих точках имеем локальный минимум


функции $\theta=\theta(u_1)$. Эти точки --- ближайшие к


$u_1$-оси, представляющие большие круги на сфере $S^3$.


Следовательно, параметр $\mu$ имеет следующий


геометрический смысл.





\begin{cor}


Для $C=-\mu^2$ минимальное расстояние от рассматриваемой поверхности


до большого круга $R (e_1 \cos u_1 + e_2 \sin u_1)$ задается 


в виде $\arcsin\mu$. 


\end{cor}





Если период функции $\theta=\theta(u_1)$ есть рациональное кратное


$2\pi$, т.\,е.


\begin{equation} \label{integral p/q}


\frac{1}{\pi} \int_{\theta_2}^{\theta_1} \sqrt{\frac{\cos^2\theta-


R^2 \sin^2\theta + \mu^2}{\cos^2\theta(R^2 \sin^2\theta-\mu^2)}}


d\theta = \frac{p}{q} \ ,


\end{equation}


где $p$, $q$ целые, то получаем замкнутую (хотя и с особенностями)


поверхность.


Интеграл в левой части соотношения \rf{integral p/q} зависит от $\mu$.


Если $\mu\rightarrow 0$, то его значение (равное половине периода


функции $\theta=\theta(u_1)$)


возрастает до бесконечности.


Поэтому существует такая последовательность


$\mu_1, \mu_2, \dotsc,\mu_k,\dotsc \rightarrow 0$, что


интеграл \rf{integral p/q} есть целое число (равное $1,2,\dotsc,


k,\dotsc$ соответственно).


Заметим, что замкнутая регулярная поверхность вращения с постоянной


отрицательной гауссовой кривизной $K=-1$ не может существовать.


Действительно, {\it не существует


замкнутой регулярной поверхности класса $C^3$ в $S^3$ с гауссовой


кривизной $K=-1$}, иначе в любой ее точке имели бы


два различных асимптотических направления.


Они никогда не совпадают, т.\,к. кручения


соответствующих асимптотических кривых


различны по знаку:


$\varkappa_1 = \sqrt{-K_e}$ и $\varkappa_2 = -\sqrt{-K_e}$, где внешняя кривизна


$K_e$ отрицательна:


$K_e = L_{11}^1 L_{22}^1 / (g_{11} g_{22}) = - 1 - R^{-2} < 0$.


По теореме Хопфа характеристика Эйлера такой поверхности


равна нулю, следовательно, рассматриваемая поверхность должна быть


гомеоморфна тору. Однако на торе не может быть


метрики постоянной отрицательной кривизны.


\medskip





{\it Случай $n=2$, $C>0$}.


В этом случае обозначим $C=\mu^2$, где $0<\mu<1$. Тогда


\begin{equation} \label{aa}


\left( \frac{d u_1}{d \theta} \right)^2 = 


\frac{\cos^2\theta - R^2 \sin^2\theta - \mu^2}{\cos^2\theta(R^2 \sin^2\theta


+\mu^2)}.


\end{equation}


Для $\theta_3$ такого, что $\sin\theta_3 = \sqrt{\frac{1-\mu^2}{1+R^2}}$,


касательная прямая к графику $u_1=u_1(\theta)$ параллельна 


$\theta$-оси (см. рис.~1c)).


Соответствующая точка порождает ребро возврата


на рассматриваемой поверхности. 


Так как $\theta_3 < \pi/2$, числитель в правой части уравнения


\rf{aa} не обращается в нуль для $0\leq \theta\leq \theta_3$.


Следовательно, график функции $\theta=\theta(u_1)$ пересекает


$u_1$-ось под ненулевым углом.


Точки пересечения соответствуют коническим точкам поверхности.


Заметим, что поверхность замкнута при любом $C>0$.


Таким образом, имеет место 





\begin{cor}


Для $n=2$ существуют три типа псевдосферических поверхностей ($K=-1$),


обобщающих тор Клиффорда $S^3$ вида \rf{gen-Cliff}


(рис.~1a), b), c)).


Для $n>2$ существует только один тип таких подмногообразий (рис.~1a)).


\end{cor}





Задача нахождения замкнутых сингулярных псевдосферических поверхностей


в $S^3$ вида \rf{gen-Cliff} является ``спектральной задачей'' в следующем 


смысле. Постоянная $C$ может рассматриваться


как спектральный параметр. Для $C>0$ имеем непрерывный спектр 


(все обобщенные торы замкнуты в этом случае) и для $C<0$


имеем последовательность дискретных ``собственных значений''


$C_1,C_2,\dotsc,C_n,\dotsc$


($\lim\limits_{n\rightarrow \infty} C_n = 0$).








\section{Заключения и открытые проблемы}





Следуя подходу, предложенному Доливой и Сантини к эволюции


кривых в $E^3$, получили спектральную задачу со спектральным


параметром для изометрических погружений $L^n$ в $E^{2n-1}$.


Для этого достаточно


рассмотреть погружения $L^n$ в сферу $S^{2n-1}$ радиуса $R$.


Радиус можно выразить в терминах параметра $\lambda$, поэтому


спектральный параметр так существенен для применения методов


теории солитонов.


В то же время вывели формулу Сыма--Тайфеля,


рассматривая предельный случай $R\rightarrow \infty$


(что означает $S^{2n-1} \rightarrow E^{2n-1}$).


Формула Сыма--Тайфеля успешно применяется в случае


погружений в $E^3$ (\cite{Sym}, \cite{Ci-FG}, \cite{Bob}),


который может рассматриваться как ``промежуточный случай''


между кривыми и многомерными многообразиями.





\begin{probl}


Вывести формулу Сыма--Тайфеля (аналогично выводу в


разделе~\ref{Symfor})


в случае поверхностей с $K=\const$, поверхностей с $H=\const$,


изотермических поверхностей и поверхностей Бианки.


\end{probl}





Специальные погружения пространств Лобачевского в евклидовы пространства 


и сферы изучались различными авторами (\cite{AG}--\cite{Ten-book}).


Мы сосредоточили внимание на замкнутых поверхностях в $S^{2n-1}$


(допускающих сингулярности), которые похожи на тор Клиффорда.


Случай $n=2$ оказывается наиболее интересным.


Обсуждение формулы \rf{integral p/q} показывает, что


для $p/q = 1$ получаем замкнутую поверхность рода~1


с единственной сингулярной линией (рис.~2a)). 


\begin{center}


\unitlength=1mm


\begin{picture}(170,68) %% width, heigth, added 5 mm for signature


\put(9,66){\special {em:graph am-ch2.bmp}} %% left X, upper Y, -, /2, +toX


\put(78,0){\mbox Рис.\,2}


\end{picture}


\end{center}








\begin{probl}


Возможно ли построить псевдосферическую поверхность 2 рода


(c отрицательной характеристикой Эйлера)


только с одной сингулярной линией (рис.~2b))?


\end{probl}





В данной работе рассмотрели псевдосферические поверхности с радиус-вектором


специального вида \rf{gen-Cliff}.


Естественно назвать ``обобщенными торами Клиффорда''


подмногообразия вида \rf{C-torus} с $c_k = c_k (u_1)$. 





\begin{probl}


Найти все обобщенные торы Клиффорда c постоянной секционной кривизной в


$S^{2n-1}$.


\end{probl}





Авторы выражают глубокую благодарность проф. Антони Сыму


за интерес, проявленный к данной работе.








\begin{thebibliography}{99}





\bibitem{ZMNP} %1


Захаров В.Е., Манаков С.В., Новиков С.П., Питаевский Л.П. 


{\it Теория солитонов. Метод обратной задачи}. -- М.: Наука, 1980. -- 319 с.





\bibitem{Sym} %2


Sym A. {\em Soliton surfaces and their application $($soliton 


geometry from spectral problems$)$} // 


Lect. Notes Phys. -- 1985. -- \No\,239. -- P.\,154--231.





\bibitem{Ci-FG} %3


Cie\'sli\'nski J.


{\em A generalized formula for integrable classes of surfaces in Lie


algebras} // J. Math. Phys. -- 1997. -- V.\,38. -- P.\,4255--4272.





\bibitem{S-1} %4


Sym A. {\it Soliton surfaces} // Lett. Nuovo Cim. -- 1982. -- V.\,33.


-- \No\,12. -- P.\,394--400.





\bibitem{S-2} %5


Sym A. {\em Soliton surfaces} II. {\em Geometric unification of solvable 


nonlinearities} // Lett. Nuovo Cim. -- 1983. -- V.\,36. -- P.\,307--312.





\bibitem{DS0} %6


Doliwa A., Santini P.


{\em The integrable dynamics of a discrete curve and the Ablowitz-Ladik


hierarchy} // 


J. Math. Phys. -- 1995. -- V.\,36. -- P.\,1259--1273.


 


\bibitem{DS2} %7


Doliwa A., Santini P.


{\em The integrable dynamics of a discrete curve} /


D.\,Levi, L.\,Vinet, P.\,Winternitz (Eds.).


Symmetries and integrability of


difference equations. -- AMS, Providence. --


CRM Proc. \& Lect. Notes. -- 1996. -- V.\,9.


-- P.\,91--102.


 





\bibitem{Am77} %8


Аминов Ю.А.


{\em О погружении областей $n$-мерного пространства Лобачевского в


$(2n-1)$-мер\-ное евклидово пространство} //


ДАН СССР. -- 1977. -- Т.\,236. -- \No\,3. -- С.\,521--524.


 


\bibitem{Am80} %9


Аминов Ю.А.


{\em Изометрические погружения областей $n$-мерного пространства Лобачевского


в $(2n-1)$-мерное евклидово пространство} // Матем. сб. -- 1980.


-- Т.\,111. -- \No\,3. -- С.\,402--433.


 


\bibitem{TT} %10


Tenenblat K., Terng C.L.


{\em B\"acklund's theorem for $n$-dimensional submanifolds of 


$R^{2n-1}$} // Ann. Math. -- 1980. -- V.\,111. -- \No\,3. -- P.\,477--490.








\bibitem{Te} %11


Terng C.L.


{\em A higher dimension generalization of the sine-Gordon equation and its


soliton theory} // 


Ann. Math. -- 1980. -- V.\,111. -- \No\,3. -- P.\,491--510.





\bibitem{ABT} %12


Ablowitz M.J., Beals R., Tenenblat K.


{\em On the solution of the generalized wave and generalized sine-Gordon


equations} //


Stud. Appl. Math. -- 1986. -- V.\,74. -- \No\,3. -- P.\,177--203.





\bibitem{FP} %13


Ferus D., Pedit F.


{\em Isometric immersions of space forms and soliton theory} //


Math. Ann. -- 1996. -- V.\,305. -- P.\,329--342.





\bibitem{CA} %14


Cie\'sli\'nski J.L., Aminov Yu.A.


{\em A geometric interpretation of the spectral problem for the generalized 


sine-Gordon system} //


J. Phys. A Math. Gen. -- 2001. -- V.\,34. -- P.\,L153--L159.





\bibitem{Mo} %15


Moore J.D.


{\em Isometric immersions of space forms in space forms} //


Pacific J. Math. -- 1972. -- V.\,40. -- \No\,1. -- P.\,157--166.





\bibitem{BT} %16


Budinich P., Trautman A. {\em The spinorial chessboard}. -- Springer-Verlag,


1988. -- 128 p.





\bibitem{Po} %17


Porteous I.R.


{\em Clifford algebras and the classical groups}. -- 


Cambridge Univ. Press, 1995. -- 295 p.





\bibitem{Ci-space} %18


Cie\'sli\'nski J.


{\em The spectral interpretation of $n$-spaces of constant negative curvature


immersed in $R^{2n-1}$} //


Phys. Lett. -- 1997. -- V.\,236. -- P.\,425--430.





\bibitem{Am88} %19


Аминов Ю.А.


{\em Изометрические погружения областей $n$-мерного пространства Лобачевского


в евклидовы пространства с плоской нормальной связностью. Модель калибровочного 


поля} // Матем. сб. -- 1988. -- Т.\,137. -- \No\,3. -- С.\,275--299.


 


\bibitem{Am98} %20


Aminov Yu.A.


{\it New ideas in differential geometry of submanifolds} //


Acta Academic, Kharkov, 1998. -- 114 p.





\bibitem{Bob} %21


Bobenko A.I. {\em Surfaces in terms of $2$ by $2$ matrices. Old and new


integrable cases} / A.P.\,Fordy and J.C.\,Wood (Eds.). Harmonic maps and 


integrable systems (Aspects of Mathematics, V.\,23),


Vieweg, Brunswick, 1994. 





\bibitem{BP} %22


Bobenko A., Pinkall U.


{\em Discrete surfaces with constant negative gaussian curvature and the


Hirota equation} // J. Diff. Geom. -- 1996. -- V.\,43. -- P.\,527--611.





\bibitem{AG} %23


Aminov Yu.A., Gontcharova O.


{\em An example of isometric immersion of a domain of $3$-dimensional 


Lobachevsky space into $E^6$ with a section as the Veronese surface} //


Math. Fiz., Analiz, Geom. -- 1999. -- V.\,6. -- P.\,3--9.





\bibitem{AR} %24


Aminov Yu.A., Rabelo M.L.


{\em On toroidal submanifolds of constant negative curvature} //


Math. Fiz., Analiz, Geom. -- 1995. -- V.\,2. -- P.\,502--513.





\bibitem{DT} %25


Dajczer M., Tojeiro R.


{\em An extension of the classical Ribaucour transformation} //


Proc. London Math. Soc. -- 2002. -- V.\,85. -- P.\,211--232.





\bibitem{Go2} %26


Gollek H.


{\em B\"acklund transformations, matrix-Riccati systems and isometric immersions


of space forms into space forms} // Coll. Math. Soc. J. Bolyai. -- 1989.


-- V.\,56. -- P.\,359--385.





\bibitem{RT} %27


Rabelo M.L., Tenenblat K.


{\em Toroidal submanifolds of constant non-positive curvature}. --


In memoriam N.I.\,Lobachevskii. -- V.\,III. -- Part~1. -- Казань,


Изд-во Казанск. ун-та, 1995. -- 


Т.\,3. -- \No\,1. -- С.\,135--159.


 


\bibitem{Ten-book} %28


Tenenblat K.


{\it Transformations of manifolds and applications to differential equations}.


-- Addison Wesley, Longman, 1998. --- 209 p.








\end{thebibliography}





\vskip 2cm





\post{\it Университет в Белостоке, институт математики,}{\qquad \it Поступила}


\post{\it Физико-технический институт низких}{\qquad 20.01.2004}


\post{\it температур им. Б.И.\,Веркина}{}


\post{\it Национальной Академии наук Украины}{}


\smallskip


\post{\it Университет в Белостоке,}{}


\post{\it институт теоретической физики}{}





%% \post{\it Uniwersytet w Bia\l ymstoku,}{}


%% \post{\it Instytut Fizyki Teoretycznej}{}








\label{end}


\end{document}




















\bibitem{Ci-dbt}


J.Cie\'sli\'nski: 


``An algebraic method to construct the Darboux matrix'', 


{\it J.Math.Phys.} {\bf 36} (1995) 5670-5706.


